% !TeX root = ../../main.tex
\subsection{Top Hodge intersections and exact vertices}
\label{subsec:s06:01}

We use Faber--Pandharipande's multipoint theorem
\cite[Theorem~1]{FP03}, together with their earlier one-point
calculation \cite[Theorem~2, (6)]{FP00}. Set $\beta_0=1$ by
convention and, for an integer $G\geq1$, define
\begin{equation}\label{eq:s06:beta-definition}
 \beta_G=\int_{\MM_{G,1}}\lambda_G\psi_1^{2G-2}.
\end{equation}
In particular $\beta_0$ is not an integral on a genus-zero one-point
space.

\begin{thm}[Top Hodge intersections]\label{thm:s06:lambda-g}
For integers $G\geq1$, $N\geq1$ and nonnegative integers
$d_1,\ldots,d_N$ with $\sum_i d_i=2G-3+N$,
\begin{equation}\label{eq:s06:lambda-g}
 \int_{\MM_{G,N}}\lambda_G\prod_{i=1}^N\psi_i^{d_i}
 =\beta_G\frac{(2G-3+N)!}{\prod_{i=1}^N d_i!}.
\end{equation}
The same formula holds for $G=0$, $N\geq3$ with $\beta_0=1$.
For a formal variable $z$,
\begin{equation}\label{eq:s06:beta-series}
 \sum_{G\geq0}\beta_Gz^{2G}=\frac{z/2}{\sin(z/2)}.
\end{equation}
As $G\to\infty$,
\begin{equation}\label{eq:s06:beta-asymptotics}
 \beta_G=2(2\pi)^{-2G}(1+O(4^{-G})),\qquad
 \frac{\beta_{G-1}}{\beta_G}=(2\pi)^2(1+O(4^{-G})).
\end{equation}
There are absolute constants $c,C>0$ such that
$c(2\pi)^{-2G}\leq\beta_G\leq C(2\pi)^{-2G}$ for every $G\geq0$.
\end{thm}
\begin{prf}
The positive-genus formula is \cite[Theorem~1]{FP03}; the
stable genus-zero formula is its equation~(1). The generating
series is \cite[Theorem~2, (6)]{FP00}, whose genus-zero value is
one. For real $s>1$, write
$\zeta(s)=\sum_{k\geq1}k^{-s}$ for the Riemann zeta function.
The Bernoulli evaluation in \cite[\S 4.2, Lemma~1]{FP00}
and \cite[(25.6.2)]{DLMF} give, for $G\geq1$,
\[
 \beta_G=(1-2^{1-2G})\frac{|B_{2G}|}{(2G)!}
       =2(1-2^{1-2G})\zeta(2G)(2\pi)^{-2G}.
\]
The Dirichlet series \cite[(25.2.1)]{DLMF} implies
$\zeta(2G)=1+O(4^{-G})$ as $G\to\infty$.
This proves both asymptotics and the uniform comparison, extending
the constants to the finitely many small genera and $G=0$.
\end{prf}

For $x\in\mathbb C$, define the entire function
\begin{equation}\label{eq:s06:phi}
 \Phi(x)=\sum_{a\geq0}\frac{(-1)^a x^a}{a!(a+1)!}
        =\frac{J_1(2\sqrt x)}{\sqrt x}.
\end{equation}
The quotient and every subsequent square-root expression in Bessel
functions mean their entire power series.

For a stable pair $(G,n)$, let $\mathbf s=(s_1,\ldots,s_n)$
be a complex vector and set
\begin{equation}\label{eq:s06:vertex-definition}
 W^{\rm top}_{G,n}(\mathbf s)
  =\int_{\MM_{G,n}}\lambda_G
                   e^{\kappa_1-\sum_{i=1}^n s_i\psi_i}.
\end{equation}
The superscript distinguishes this vertex integral from other
intersection generating functions. The local coefficient order
$m=2G-3+n$ is nonnegative for every stable type.

The whole exponential specialization
\eqref{eq:s02:kappa-exponential} is especially useful here. It
replaces the background by new cotangent insertions, after which
\eqref{eq:s06:lambda-g} determines every coefficient.

\begin{prp}[Exact vertices]\label{prp:s06:exact-vertex}
For every stable pair $(G,n)$, with $m=2G-3+n$,
\begin{equation}\label{eq:s06:exact-vertex}
 W^{\rm top}_{G,n}(\mathbf s)
 =\beta_G m![x^m]\Phi(x)^{-m-1}
                       \prod_{i=1}^n J_0(2\sqrt{s_i x}).
\end{equation}
\end{prp}
\begin{prf}
Apply \eqref{eq:s02:kappa-exponential} with scalar parameter one.
A term with $k\geq0$ forgotten points has coefficient $1/k!$ and
new-point factor $\psi(1-e^{-\psi})$ at each point.
The Hodge insertion pulls back unchanged. The new powers have
minimum degree two, so Lemma~\ref{lmm:s02:rational-tail-degree}
replaces the old cotangent pullbacks by their intrinsic classes.
This is the inversion of the entire kappa background; there is no
further background exponential to pull back.

The multinomial formula contributes $\beta_G(m+k)!$ and a
factorial denominator for each descendant degree. Thus the old
marking $i$ contributes
\[
 \sum_{a\geq0}\frac{(-s_i x)^a}{(a!)^2}
       =J_0(2\sqrt{s_i x}),
\]
whereas a new marking contributes
\[
 \sum_{a\geq1}\frac{(-1)^{a-1}x^{a+1}}{a!(a+1)!}
       =x(1-\Phi(x)).
\]
After removing $x^k$ from the coefficient extraction, the vertex
integral is therefore
\[
 \beta_G[x^m]\prod_{i=1}^nJ_0(2\sqrt{s_i x})
       \sum_{k\geq0}\frac{(m+k)!}{k!}(1-\Phi(x))^k.
\]
This sum is coefficientwise finite because $1-\Phi(x)$ has zero
constant term. The binomial series makes it
$m!\Phi(x)^{-m-1}$, proving the formula.

For an unmarked vertex $G\geq2$, one has $m\geq1$. Its $k=0$
term vanishes by degree, and every nonzero term has at least one
marking, as required by the positive-genus multipoint theorem.
For $G=0$ the original stability already gives $n\geq3$.
These observations cover the ranges excluded from a literal
application of \eqref{eq:s06:lambda-g} with $N=0$.
\end{prf}

\begin{crl}[Numerical genus-zero reduction]\label{crl:s06:genus-zero}
For every stable type,
\begin{equation}\label{eq:s06:genus-zero}
 W^{\rm top}_{G,n}(\mathbf s)
 =\beta_G W^{\rm top}_{0,2G+n}(\mathbf s,0,\ldots,0).
\end{equation}
\end{crl}
\begin{prf}
The type on the right is stable since $2G+n=m+3\geq3$.
It has the same coefficient order $m$, and its additional $2G$
old-marking factors are $J_0(0)=1$ in
\eqref{eq:s06:exact-vertex}. Its one-point convention is $\beta_0=1$.
\end{prf}

\begin{rmk}[Classical inverse]
This is an equality of numerical generating functions on different
moduli spaces. The inverse of $x\mapsto x\Phi(x)$ is the classical
genus-zero kappa series of \cite[(0.8)]{KMZ96v1}, with its earlier
Zograf attribution. For a formal variable $t$ and an integer $N\geq3$,
Lagrange inversion and \eqref{eq:s06:exact-vertex} give its
$t^{N-2}$ coefficient as
\[
 \frac{W^{\rm top}_{0,N}(\mathbf0)}{(N-2)!}
 =\frac{\displaystyle\int_{\MM_{0,N}}\kappa_1^{N-3}}
        {(N-2)!(N-3)!}.
\]
This makes the conversion from raw kappa powers to exponential
integrals explicit. The inverse and critical behavior are also
analyzed in \cite[\S 6, proof of 6.3]{MaZ00}. The uniform vertex
and complete tree estimates needed here will be proved directly.
\end{rmk}
